\documentclass[12pt]{article}
\usepackage[utf8]{inputenc}
\usepackage[margin=1in]{geometry}
\usepackage{amsmath}
\usepackage{graphicx}
\usepackage{booktabs}
\usepackage{hyperref}
\usepackage[round]{natbib}

\title{Superficial Superior Colliculus Is Causally Involved in Complex Figure Detection and Encodes Figure-Ground Signals in Mice}
\author{Author A$^{1}$, Author B$^{1,2}$, Author C$^{1}$ \\
\small $^{1}$Department of Neurophysiology, Donders Institute, Radboud University, Nijmegen, Netherlands \\
\small $^{2}$Neuroscience Program, Donders Graduate School, Nijmegen, Netherlands}
\date{}

\begin{document}
\maketitle

\begin{abstract}
Object detection on complex backgrounds is typically attributed to visual cortex, yet the superior colliculus (SC) receives direct retinal input and can support visual behaviors after cortical damage. Whether the superficial SC (sSC) contributes to detecting figures defined by complex features on textured backgrounds has not been tested. Here, we used bilateral optogenetic inhibition of sSC in GAD2-Cre mice ($N = 8$) performing three figure detection tasks (contrast, orientation, and phase differences) and found that sSC inhibition significantly decreased performance on all three tasks, with the strongest effects at early latencies (0--100 ms post-stimulus). Electrophysiological recordings in behaving C57BL/6J mice ($N = 5$) revealed that sSC neurons respond more strongly to figure elements than identical ground elements for contrast and orientation tasks. A linear SVM decoder discriminated figure from ground above chance for all three tasks, including phase, where population-level response differences were not significant. Critically, figure-ground discriminability ($d'$) was significantly higher on correct trials than error trials, linking sSC neural activity to behavioral decisions. A small group of putative multisensory neurons at slightly deeper locations showed task-related responses peaking at approximately 700 ms, suggesting integration of visual and decision-related signals. These findings establish that the sSC is causally involved in complex figure detection and encodes behaviorally relevant figure-ground information.
\end{abstract}

\section{Introduction}

Detecting objects against complex backgrounds is a fundamental visual computation traditionally attributed to cortical circuits, particularly primary visual cortex (V1) and its feedback connections with higher visual areas \citep{lamme1995figure}. Figure-ground modulation---the enhanced neural response to figure elements compared to identical stimulus elements when they form part of the background---has been extensively characterized in V1 of primates and is thought to depend on recurrent cortical processing and top-down feedback \citep{lamme1995figure}.

However, the superior colliculus (SC) receives direct retinal input and plays an established role in visual orienting, target selection, and simple stimulus detection \citep{schneider1969sc, dean1989sc}. In mice, the SC is the largest recipient of retinal ganglion cell projections, and behavioral studies have shown that mice can still detect simple visual stimuli after V1 silencing or ablation, with this residual capacity likely mediated by the SC \citep{feinberg2021v1, zhao2014sc}. Recent work has demonstrated that mouse sSC encodes visual features of complex objects independently of V1 input \citep{gale2014sc, wang2010sc}, but whether this encoding supports detection of figures on complex, non-homogeneous backgrounds---as opposed to detecting isolated stimuli on uniform backgrounds---has not been investigated.

The critical question is whether the sSC contributes causally to figure detection when figures are defined by differences in contrast, orientation, or phase relative to a textured background. These feature differences require processing beyond simple luminance change detection and are typically associated with cortical mechanisms \citep{lamme1995figure}. If the sSC is causally involved in such tasks, it would significantly expand our understanding of the neural circuits supporting object detection beyond the cortical figure-ground modulation framework.

Here, we address this question using bilateral optogenetic inhibition of sSC combined with electrophysiological recordings during figure detection behavior in mice. We employ three task variants that require detection of figures differing from background by contrast, orientation, or phase, and we test whether sSC neural activity encodes figure-ground information in a behaviorally relevant manner.

\section{Methods}

\subsection{Animals}

Three groups of mice were used (Table~\ref{tab:animals}). All mice were housed on a reversed 12-hour light-dark cycle with ad libitum food access, and experiments were conducted during the dark cycle.

\begin{table}[htbp]
\centering
\caption{Animals used in each experiment.}
\label{tab:animals}
\begin{tabular}{llccc}
\toprule
Experiment & Strain & $N$ & Sex & Age \\
\midrule
Behaving electrophysiology & C57BL/6J & 5 & All male & 2--5 months \\
Behavior + optogenetics & GAD2-Cre & 8 & 6M, 2F & 2--5 months \\
Anesthetized electrophysiology & GAD2-Cre & 3 & All male & 2--5 months \\
\bottomrule
\end{tabular}
\end{table}

\subsection{Viral Injection and Optogenetic Inhibition}

GAD2-Cre mice received bilateral injections of Cre-dependent ChR2 (ssAAV-9/2-hEF1a-dlox-hChR2(H134R)\_mCherry, titer $5.4 \times 10^{12}$) into the sSC, followed by optic fiber implantation. Blue laser light activated inhibitory (GABAergic) neurons, suppressing overall sSC activity. Anesthetized control recordings confirmed that optogenetic inhibition reduced visually evoked responses by 33\% on average ($p < 0.001$), though silencing was not complete. A blue LED above the mouse head flashed at random intervals to prevent behavioral cueing by the laser.

\subsection{Behavioral Task}

Mice performed a two-alternative forced choice (2-AFC) figure detection task. Figures were presented on the left or right side of a textured background, and mice reported figure location by licking the corresponding side of a Y-shaped spout. Three task variants used figures differing from background by contrast, orientation, or phase. Response was allowed from 200 ms after stimulus onset. Mice completed 100--250 trials per session over 2--5 months of training.

Optogenetic inhibition was applied at different latencies (0--200 ms) after stimulus onset to assess the temporal window of sSC involvement.

\subsection{Electrophysiology}

Silicon probes (Neuropixels or 32-channel linear arrays) targeted the sSC. Neurons were classified by receptive field (RF) location relative to the figure: inside the figure, on the figure edge, or outside the figure. Population responses were compared between figure and ground conditions using cluster-based permutation tests. A linear SVM classifier with a 50-ms sliding window (10-ms steps) decoded figure versus ground identity. Behavioral relevance was assessed by comparing figure-ground discriminability ($d'$ from ROC analysis) between hit and error trials.

\section{Results}

\subsection{Optogenetic Inhibition of sSC Impairs Figure Detection}

Bilateral inhibition of sSC significantly decreased task accuracy for all three figure detection tasks---contrast, orientation, and phase (Table~\ref{tab:opto}). The temporal profile of the inhibition effect was informative: the strongest impairment occurred when inhibition was applied at early latencies (0--100 ms after stimulus onset), with diminished effects at later latencies (100--200 ms). This temporal profile is consistent with a role for sSC in the initial feedforward processing of visual information that supports figure detection.

\begin{table}[htbp]
\centering
\caption{Temporal profile of optogenetic inhibition effects on figure detection accuracy.}
\label{tab:opto}
\begin{tabular}{lccc}
\toprule
Inhibition Latency (ms) & Contrast & Orientation & Phase \\
\midrule
0--100 (early) & Strongest effect & Strongest effect & Strongest effect \\
100--200 (late) & Diminished & Diminished & Diminished \\
No laser (baseline) & Normal & Normal & Normal \\
\bottomrule
\end{tabular}
\end{table}

\subsection{sSC Neurons Encode Figure-Ground Information}

Electrophysiological recordings from neurons with RFs inside the figure region revealed that population responses were significantly stronger when the RF was in the figure compared to when the identical stimulus element was part of the background. This figure-ground modulation was statistically significant for contrast and orientation tasks (cluster-based permutation tests, $p < 0.05$), while the phase task did not reach significance at the population level.

Neurons with RFs on the figure edge also showed figure-ground modulation for contrast and orientation tasks, with significant time clusters where figure responses exceeded ground responses.

\subsection{SVM Decoding Reveals Distributed Figure-Ground Information}

A linear SVM decoder applied to population activity from RF-edge neurons could discriminate figure from ground above chance for all three tasks, including the phase task where population response differences were not significant (Table~\ref{tab:svm}). This indicates that figure-ground information for phase stimuli is present at the single-neuron level in consistent but subtle response differences that are detectable by a classifier but do not reach significance in the population average.

\begin{table}[htbp]
\centering
\caption{SVM decoding performance for figure versus ground discrimination.}
\label{tab:svm}
\begin{tabular}{lccc}
\toprule
Task & Above Chance & Significant Time Bins & Population Response \\
\midrule
Contrast & Yes & Multiple & Significant \\
Orientation & Yes & Multiple & Significant \\
Phase & Yes & Multiple & Not significant \\
\bottomrule
\end{tabular}
\end{table}

\subsection{Neural Discriminability Predicts Behavioral Success}

The most important finding linking sSC activity to behavior was that figure-ground discriminability ($d'$) was significantly higher on correct trials (hits) than on error trials for all three tasks. Because the stimulus in the RF was physically identical between hits and errors (only the mouse's behavioral response differed), the difference in neural discriminability must reflect trial-to-trial variability in the neural representation that directly influences the behavioral decision. This establishes a functional link between sSC figure-ground coding and perceptual performance.

\subsection{Putative Multisensory Neurons at Decision Time}

A small group of neurons recorded at slightly deeper locations than the sSC showed task-related responses with a distinctive temporal profile: activity peaked at approximately 700 ms after stimulus onset, coinciding with the time between lick spout presentation and the licking response. These putative multisensory neurons showed figure-ground modulation similar to the visual neurons but at a much later latency, suggesting they may integrate visual information with motor or decision-related signals at the time of behavioral report.

\section{Discussion}

This study provides causal and correlational evidence that the superficial superior colliculus contributes to complex figure detection in mice. The finding that bilateral sSC inhibition impairs detection of figures defined by contrast, orientation, and phase differences extends the known role of the SC from simple stimulus detection to figure-ground segregation on complex textured backgrounds \citep{schneider1969sc, dean1989sc, boehnke2011sc}.

\subsection{sSC Involvement Beyond Simple Detection}

The SC has traditionally been associated with detecting and orienting toward salient stimuli in the environment \citep{schneider1969sc}. Our results demonstrate that sSC processes contribute to a more sophisticated computation: distinguishing a figure from a background when both consist of textured elements differing only in a specific feature dimension. This is particularly notable for the phase task, where the luminance and contrast of figure and ground elements are matched, requiring detection based on spatial phase alignment. The finding that sSC neurons carry decodable phase-defined figure-ground information, even when population averages do not reach significance, suggests that the sSC may contribute to perceptual decisions based on subtle stimulus differences \citep{gale2014sc, resulaj2018svm}.

\subsection{Temporal Window of sSC Contribution}

The finding that optogenetic inhibition was most effective in the 0--100 ms window after stimulus onset suggests that the sSC's contribution is primarily during the initial feedforward sweep of visual processing. This is consistent with the known short-latency visual responses in the sSC and suggests that the SC provides an early figure detection signal that may be used by downstream decision circuits \citep{zhao2014sc, wang2010sc}.

\subsection{Relationship to Cortical Figure-Ground Processing}

Figure-ground modulation has been extensively studied in V1, where it is thought to depend on recurrent processing and feedback from higher visual areas \citep{lamme1995figure}. Our finding that the sSC also encodes figure-ground information raises the question of whether these represent independent parallel pathways or interact. The SC receives both direct retinal input and cortical projections from V1, and V1 in turn receives SC input via the thalamus \citep{zingg2017sc}. Future experiments combining cortical silencing with sSC recordings will be needed to determine the extent to which sSC figure-ground signals depend on cortical feedback.

\subsection{Limitations}

The partial (33\%) reduction in sSC activity achieved by optogenetic inhibition means that our behavioral effects likely underestimate the full contribution of sSC to figure detection. The relatively small electrophysiology sample ($N = 5$ mice) limits the power to detect subtle differences, particularly for the phase task. The identity of the putative multisensory neurons and their precise role in the decision process requires further characterization.

\section{Conclusion}

The superficial superior colliculus is causally involved in detecting figures defined by complex feature differences on textured backgrounds and encodes figure-ground information that predicts behavioral success. These findings expand the role of the SC beyond simple stimulus detection to include perceptual computations previously attributed primarily to cortical circuits, and identify the sSC as an important component of the neural circuitry for object detection.

\bibliographystyle{plainnat}
\bibliography{references}

\end{document}
